\section{Discussion}

The confirmatory analysis \PrimaryInterpretation. Against the preselected
PGNN-style comparator, the lower bootstrap endpoint clears the 5\% practical
margin and the exact seed-level test rejects that margin. The classical
baseline refines the interpretation of this result. A gain-0.5
context filter followed by the nominal transition reaches 0.111 NRMSE-25,
compared with 0.119 for CFR-WM and 0.132 for PGNN-style. The confirmatory H1
remains supported for its prespecified comparison. The component analysis
attributes the aggregate advantage to forecast-origin filtering rather than to
the learned residual transition.

The fixed-weight scale analysis strengthens that attribution. Removing the
residual at inference gives the lowest mean in clean ID, corrupt ID, corrupt
OOD, and the aggregate; increasing the multiplier from 0 to 0.5 monotonically
increases aggregate NRMSE-25. With transition weights held fixed, the fitted
correction increases the reported cell-balanced point error. The nominal prior
and context correction capture the useful structure in this benchmark.

The regime decomposition explains when the filter helps relative to PGNN-style.
The corrupted and shifted regimes improve by \CorruptImprovement\%, whereas
the clean regime changes by \CleanImprovement\%; consequently H2
\RobustnessInterpretation. The synchronous two-step outage produces only a
2.7\% increase over clean-context CFR-WM, versus 51.9\% for PGNN-style. This
extends the missing-at-random evidence to a short, synchronized loss of all
channels and indicates that nominal context propagation is valuable during
brief sensor interruptions.

Equation~\eqref{eq:optimal-gain} turns the gain trade-off into an estimator
design rule. Equal trust is optimal only for equal prediction and measurement
error variances with zero cross-covariance. Per-channel estimates of
$v_p$, $v_\epsilon$, and $c$ from training or validation innovations would
produce fixed channel gains without using test outcomes; a causal network could
instead condition them on current sensor metadata. Either construction matches
the observed heterogeneity better than one scalar gain. The pendulum audit also
shows why gain semantics and state constraints must be separated: 75.4\% of
missing sine/cosine entries change under the subsequent unit-circle projection,
even though their direct measurement correction is exactly zero.

The measured-flight evaluation exposes a state-component and uncertainty
trade-off. The matched unfiltered residual model has lower mean full-state
NRMSE in every condition--horizon comparison, whereas CFR-WM has lower 50-step
position error in all four conditions. On recorded contexts, velocity error
rises from \ExternalUnfilteredRecordedLongVelocityMean{} to
\ExternalCFRRecordedLongVelocityMean{}\,m\,s$^{-1}$. At the same horizon,
CFR-WM's nominal 90\% marginal coverage is only 0.524--0.667 and is below the
unfiltered model in all four conditions. Narrower intervals accompany higher
NLL and calibration error, consistent with missing cross-state dependence and
unrepresented model-form error in the diagonal Gaussian forecast.

These results connect estimator design to the control objective. Uniform NLL
weights every normalized state component equally, while a controller can value
position, velocity, or constraint-related coordinates differently.
Equation~\eqref{eq:weighted-nll} makes such weights explicit and requires them
to be fixed from the downstream specification rather than selected from test
performance. Correlated covariance or calibrated prediction sets would also be
needed when constraints depend jointly on several state coordinates.

The measured runtime places CFR-WM at
\CFRLatencyMean{}\,ms per trajectory for an eight-step context and 25-step
rollout on the stated GPU and represents throughput for 96 simultaneous
trajectories. Online control requires separate timing and closed-loop tests for
a specified controller, plant, sampling deadline, and disturbance model
\cite{hewing2020_learning_mpc}.
